\lyseis{Α3.7}{Τυποποιημένη μορφή μεγάλων αριθμών}
\sol{1} α) $1{,}5\cdot10^{5}$ \quad β) $2{,}1\cdot10^{4}$ \quad γ) $1{,}3\cdot10^{6}$ \quad δ) $2{,}03\cdot10^{6}$.
\sol{2} α) $23\,100$ \quad β) $1\,360\,000\,000$ \quad γ) $137$ \quad δ) $1\,000\,000$.
\sol{3} α) $2\,000>300$, άρα $>$. β) $0{,}1\cdot100=10$, άρα $=$. γ) $17\,000>3\,000$, άρα $>$. δ) $600\,000=0{,}6\cdot1\,000\,000=600\,000$, άρα $=$.
\sol{4} Σε τυποποιημένη μορφή είναι οι $4{,}5\cdot10^{3}$, $9{,}99\cdot10^{8}$ και $1\cdot10^{6}$. Οι $45\cdot10^2$ και $10\cdot10^5$ δεν είναι, γιατί ο πρώτος παράγοντας δεν είναι μικρότερος του $10$. Ο $0{,}45\cdot10^4$ δεν είναι, γιατί το ακέραιο μέρος είναι $0$.
\sol{5} $45\cdot10^{2}=4\,500=4{,}5\cdot10^{3}$. $0{,}45\cdot10^{4}=4\,500=4{,}5\cdot10^{3}$. $10\cdot10^{5}=1\,000\,000=1\cdot10^{6}$.
\sol{6} α) $7\cdot10^{3}$ \quad β) $4{,}8\cdot10^{8}$ \quad γ) $1{,}25\cdot10^{10}$ \quad δ) $6{,}021\cdot10^{6}$.
\sol{7} α) Σωστό. β) Λάθος: $2{,}5\cdot10^{4}=25\,000$. γ) Σωστό. δ) Σωστό. ε) Λάθος: $3\cdot10^{2}=300$. στ) Σωστό.
\sol{8} α) $3{,}844\cdot10^{5}$ \quad β) $1{,}496\cdot10^{8}$ \quad γ) $4{,}5\cdot10^{9}$ \quad δ) $1{,}36\cdot10^{10}$ \quad ε) $5{,}9\cdot10^{9}$.
\sol{9} Ένα έτος έχει $365\cdot24\cdot60=525\,600$ λεπτά. Τα $150$ έτη έχουν $150\cdot525\,600=78\,840\,000$ λεπτά, δηλαδή $7{,}884\cdot10^{7}$ λεπτά.
\sol{10} Ερμής ($5{,}8\cdot10^{7}$), Άρης ($2{,}3\cdot10^{8}$), Δίας ($7{,}8\cdot10^{8}$), Ποσειδώνας ($4{,}5\cdot10^{9}$). Συγκρίνουμε πρώτα τους εκθέτες και, αν είναι ίσοι, τους $\alpha$.
\sol{11} $3{,}2\mathrm{E}8=3{,}2\cdot10^{8}=320\,000\,000$ και $4{,}05\mathrm{E}6=4{,}05\cdot10^{6}=4\,050\,000$.
\sol{12} Κίνα: $1{,}4\cdot10^{9}$. Ελλάδα: $1{,}04\cdot10^{7}$. Μεγαλύτερος είναι ο πληθυσμός της Κίνας, γιατί έχει μεγαλύτερο εκθέτη ($9>7$).
\sol{13} Ο πρώτος παράγοντας ($12$) δεν είναι μικρότερος του $10$. Σωστά: $12\cdot10^{5}=1\,200\,000=1{,}2\cdot10^{6}$.
\sol{14} Έχει $\nu+1$ ψηφία. Οι $4\cdot10^{9}=4\,000\,000\,000$ και $4{,}5\cdot10^{9}=4\,500\,000\,000$ έχουν και οι δύο $10$ ψηφία.
\sol{15} Ο $5\,400\,000$ έχει $7$ ψηφία, άρα η υποδιαστολή πρέπει να μετακινηθεί $6$ θέσεις: $5{,}4\cdot10^{6}$. Στο $5{,}4\cdot10^{5}$ έχουμε μόνο $540\,000$.
\sol{16} Βρίσκεται ανάμεσα στο $10^{6}=1\,000\,000$ και στο $10^{7}=10\,000\,000$, γιατί $4{,}5\cdot10^{6}=4\,500\,000$.
\sol{17} Ελεύθερη απάντηση. Π.χ. πληθυσμός της Γης $8\,000\,000\,000=8\cdot10^{9}$. Κάθε αριθμός πρέπει να έχει πρώτο παράγοντα από $1$ ως κάτω από $10$.
\sol{18} $365\cdot24\cdot3\,600=31\,536\,000=3{,}1536\cdot10^{7}$ δευτερόλεπτα.
